\documentclass{article}
\usepackage{graphicx}
\usepackage{geometry}
\usepackage{amsmath}
\usepackage{booktabs}

\begin{document}
\section*{NOTES TO UNAUDITED CONDENSED CONSOLIDATED FINANCIAL STATEMENTS}

\subsection*{1. BASIS OF PRESENTATION AND SUMMARY OF SIGNIFICANT ACCOUNTING POLICIES}

\subsubsection*{A. BASIS OF PRESENTATION}

The unaudited interim condensed consolidated financial statements of RLI Corp. (the Company) and subsidiaries have been prepared in accordance with generally accepted accounting principles in the United States of America (GAAP). These condensed consolidated financial statements do not include all the disclosures required by GAAP for annual financial statements and should be read in conjunction with our 2023 Annual Report on Form 10-K. In the opinion of the Company's management, the condensed consolidated financial statements reflect all adjustments, of a normal and recurring nature, that are necessary for fair financial statement presentation. The results of operations for any interim period are not necessarily indicative of the operating results for a full year.

The preparation of the unaudited condensed consolidated financial statements requires management to make estimates and assumptions relating to the reported amounts of assets and liabilities, the disclosure of contingent assets and liabilities at the date of the unaudited condensed consolidated financial statements and the reported amounts of revenue and expenses during the period. These estimates are inherently subject to change and actual results could differ significantly from these estimates.

\subsubsection*{B. ADOPTED ACCOUNTING STANDARDS}

No new accounting standards applicable in 2024 materially impact our financial statements.

\subsubsection*{C. PROSPECTIVE ACCOUNTING STANDARDS}

2023-07—Segment Reporting (Topic 280): Improvements to Reportable Segment Disclosures
The guidance in ASU 2023-07 was designed to improve reportable segment disclosure requirements, primarily through enhanced disclosures about significant segment expenses. This ASU is effective for fiscal years beginning after December 15, 2023, and interim periods within fiscal years beginning after December 15, 2024. Although the Company continues to evaluate the impact of adopting this new accounting standard, the amendments are disclosure-related and should not have a material impact on our financial statements.

2023-09—Income Taxes (Topic 740): Improvements to Income Tax Disclosures
The guidance in ASU 2023-09 was designed to increase transparency about income tax information through improvements to the rate reconciliation and disclosure of income taxes. This ASU is effective for fiscal years beginning after December 15, 2024. Although the Company continues to evaluate the impact of adopting this new accounting standard, the amendments are disclosure-related and should not have a material impact on our financial statements.

\subsubsection*{D. REINSURANCE}
Ceded unearned premiums and reinsurance balances recoverable on unpaid losses and settlement expenses are reported separately as an asset, rather than being netted with the related liability, since reinsurance does not relieve the Company of our liability to policyholders. Such balances are subject to the credit risk associated with the individual reinsurer. We continually monitor the financial condition of our reinsurers and actively follow up on any past due or disputed amounts. As part of our monitoring efforts, we review reinsurers' annual financial statements and Securities and Exchange Commission filings for those that are publicly traded. We also review insurance industry developments that may impact the financial condition of our reinsurers. We analyze the credit risk associated with our reinsurance balances recoverable by monitoring the A.M. Best and Standard \& Poor's (S\&P) ratings of our reinsurers. We subject our reinsurance balances recoverable to detailed recoverability tests, including a segment-based analysis using the average default rating percentage by S\&P rating, which assists the Company in assessing the sufficiency of its allowance. Additionally, we perform an in-depth reinsurer financial condition analysis prior to the renewal of our reinsurance placements.

Our policy is to charge to earnings, in the form of an allowance, an estimate of unrecoverable amounts from reinsurers. This allowance is reviewed on an ongoing basis to ensure that the amount makes a reasonable provision for reinsurance balances that we may be unable to recover. Once regulatory action (such as receivership, finding of insolvency, order of conservation or order of liquidation is taken against a reinsurer, the paid and unpaid recoverable balances for the reinsurer are specifically identified and written off through either a cut or allowance for estimated unrecoverable amounts from reinsurers. When we write off such a balance, it is done in full. We then re-evaluate the overall allowance and determine whether the balance is sufficient and, if needed, an additional allowance is recognized.

The allowances for uncollectible amounts on paid and unpaid reinsurance balances recoverable were \$16 million and \$10 million, respectively, at September 30, 2024. At December 31, 2023, the amounts were \$16 million and \$11 million, respectively. Changes in the allowances were due to changes in the amount of reinsurance balances outstanding, the composition of reinsurers from whom the balances were recoverable and their associated S\&P default ratings. No write-offs were applied to the allowances in the first nine months of 2024 and less than \$1 million was recovered.

\subsubsection*{E. INTANGIBLE ASSETS}
The composition of goodwill and intangible assets at September 30, 2024 and December 31, 2023 is detailed in the following table:

\begin{table}[h!]
\centering
\begin{tabular}{lcc}
& September 30, 2024 & December 31, 2023 \\
\hline
Goodwill & \$ 43,416 & \$ 43,416 \\
\hline
Surety & \$ 5,245 & \$ 5,245 \\
Casualty & \$ 456 & \$ 456 \\
Total goodwill & \$ 5,701 & \$ 5,701 \\
\hline
Indefinite-lived intangibles & \$ 750 & \$ 750 \\
Total goodwill and intangibles & \$ 37,562 & \$ 37,562 \\
\end{tabular}
\end{table}

Annual impairment assessments were performed on our goodwill and state insurance license indefinite-lived intangible assets during the second quarter of 2024. Based upon these reviews, none of the assets were impaired. In addition, there were no triggering events as of September 30, 2024 that would suggest an updated impairment test would be needed for our goodwill and intangible assets.

\subsubsection*{F. EARNINGS PER SHARE}
Basic earnings per share (EPS) is computed by dividing income available to common shareholders by the weighted-average number of common shares outstanding for the period. Diluted EPS reflects the dilution that could occur if securities or other contracts to issue common stock or common stock equivalents were exercised or converted into common stock. When inclusion of these items increases the earnings per share or reduces the loss per share, the effect on earnings is anti-dilutive. Under these circumstances, the diluted net earnings or net loss per share is computed excluding these items. The following represents a reconciliation of the numerator and denominator of the basic and diluted EPS computations contained in the unaudited condensed consolidated financial statements:


\end{document}